% META: {"id": "TSPE-COMBINATOIRE-QCM", "chapitre": "TSPE-COMBINATOIRE", "type_objet": "qcm", "genere_depuis": "chapitres/TSPE-COMBINATOIRE/qcm/TSPE-COMBINATOIRE-QCM.json", "status": "generated"}
% Fichier genere par scripts/build_qcm_tex.py — ne pas editer a la main.

\section*{\textcolor{chapcolor}{\MakeUppercase{Combinatoire et denombrement}}}

\begin{center}
\textit{Pour chaque question, une seule reponse est exacte.}
\end{center}

\begin{enumerate}

\item[] \textbf{Capacite C2}

\item \textbf{[Q1]} Le nombre de k-listes d'un ensemble a n elements est :
  \begin{enumerate}[label=\Alph*.]
    \item $n^k$
    \item $k^n$
    \item $n!$
    \item $\binom{n}{k}$
  \end{enumerate}

\item \textbf{[Q2]} Le nombre d'arrangements sans repetition $A_n^k$ est :
  \begin{enumerate}[label=\Alph*.]
    \item $n^k$
    \item $\frac{n!}{k!}$
    \item $\frac{n!}{(n-k)!}$
    \item $\binom{n}{k}$
  \end{enumerate}

\item[] \textbf{Capacite C1}

\item \textbf{[Q3]} Un choix se fait en trois etapes offrant successivement $2$, $3$ puis $4$ possibilites. Quelle demarche represente correctement le denombrement ?
  \begin{enumerate}[label=\Alph*.]
    \item Additionner les branches : $2+3+4=9$
    \item Utiliser un arbre et le principe multiplicatif : $2\times3\times4=24$
    \item Choisir trois elements parmi quatre : $\binom{4}{3}=4$
    \item Permuter les quatre possibilites finales : $4!=24$
  \end{enumerate}

\item[] \textbf{Capacite C3}

\item \textbf{[Q4]} Quelle egalite traduit le denombrement de toutes les parties d'un ensemble a $n$ elements selon leur cardinal ?
  \begin{enumerate}[label=\Alph*.]
    \item $\displaystyle\sum_{k=0}^{n} \binom{n}{k}=2^n$
    \item $\displaystyle\sum_{k=0}^{n} \binom{n}{k}=n$
    \item $\displaystyle\prod_{k=0}^{n} \binom{n}{k}=2^n$
    \item $\displaystyle\sum_{k=1}^{n} k=n!$
  \end{enumerate}

\item[] \textbf{Capacite C4}

\item \textbf{[Q5]} Le triangle de Pascal donne $\binom{n}{k} + \binom{n}{k+1} =$
  \begin{enumerate}[label=\Alph*.]
    \item $\binom{n}{k+2}$
    \item $\binom{n+2}{k}$
    \item $\binom{n+1}{k}$
    \item $\binom{n+1}{k+1}$
  \end{enumerate}

\end{enumerate}

% NEXUS-QCM-TEACHER-ONLY-BEGIN
\ifnxVersionProfesseur
\clearpage
\section*{Cle de correction — reservee au professeur}

\begin{center}
\begin{tabular}{lll}
\hline
Question & Capacite & Reponse exacte \\
\hline
Q1 & C2 & \textbf{A} \\
Q2 & C2 & \textbf{C} \\
Q3 & C1 & \textbf{B} \\
Q4 & C3 & \textbf{A} \\
Q5 & C4 & \textbf{D} \\
\hline
\end{tabular}
\end{center}

\subsection*{Diagnostic des reponses erronees}

\begin{itemize}
  \item \textbf{Q1}
  \begin{itemize}
    \item \textbf{B} — Inversion de la base et de l'exposant. \emph{Renvoi : C2.}
    \item \textbf{C} — Formule des permutations. \emph{Renvoi : C2.}
    \item \textbf{D} — Formule des combinaisons. \emph{Renvoi : C2.}
  \end{itemize}
  \item \textbf{Q2}
  \begin{itemize}
    \item \textbf{A} — Formule avec remise. \emph{Renvoi : C2.}
    \item \textbf{B} — Le denominateur $k!$ ne retire pas les $n-k$ facteurs inutilises ; une k-liste sans repetition conserve exactement $n!/(n-k)!$. \emph{Renvoi : C2.}
    \item \textbf{D} — Formule sans ordre. \emph{Renvoi : C2.}
  \end{itemize}
  \item \textbf{Q3}
  \begin{itemize}
    \item \textbf{A} — Le principe additif convient a des cas disjoints alternatifs ; ici chaque choix du premier niveau se combine avec les choix suivants. \emph{Renvoi : C1.}
    \item \textbf{C} — La combinaison $\binom{4}{3}$ choisirait un sous-ensemble sans ordre ; elle ne represente pas trois etapes successives de tailles differentes. \emph{Renvoi : C1.}
    \item \textbf{D} — Le nombre $4!$ correspond aux permutations de quatre objets distincts ; son egalite numerique a $24$ est fortuite et le modele est inadapte. \emph{Renvoi : C1.}
  \end{itemize}
  \item \textbf{Q4}
  \begin{itemize}
    \item \textbf{B} — La somme compte toutes les parties, pas seulement les $n$ singletons ; un ensemble a $n$ elements possede $2^n$ parties. \emph{Renvoi : C3.}
    \item \textbf{C} — Les classes de parties de cardinal $k$ sont disjointes et doivent etre additionnees ; leur produit ne denombre pas leur reunion. \emph{Renvoi : C3.}
    \item \textbf{D} — La somme $1+\cdots+n$ ne fait intervenir aucun coefficient binomial et vaut $n(n+1)/2$, pas $n!$ en general. \emph{Renvoi : C3.}
  \end{itemize}
  \item \textbf{Q5}
  \begin{itemize}
    \item \textbf{A} — Le choix $\binom{n}{k+2}$ decale deux fois l'indice inferieur sans passer a la ligne $n+1$ du triangle de Pascal. \emph{Renvoi : C4.}
    \item \textbf{B} — Le choix $\binom{n+2}{k}$ avance de deux lignes alors que la somme de deux coefficients voisins produit un coefficient de la seule ligne $n+1$. \emph{Renvoi : C4.}
    \item \textbf{C} — Pour obtenir $\binom{n+1}{k}$, il faudrait additionner $\binom{n}{k-1}$ et $\binom{n}{k}$ ; ici les indices sont $k$ et $k+1$. \emph{Renvoi : C4.}
  \end{itemize}
\end{itemize}
\fi
% NEXUS-QCM-TEACHER-ONLY-END
